\documentclass[11pt,a4paper]{article}
\usepackage[margin=2.5cm]{geometry}
\usepackage{amsmath}
\usepackage{booktabs}
\usepackage{listings}
\usepackage{xcolor}

\lstset{
  basicstyle=\ttfamily\small,
  frame=single,
  backgroundcolor=\color{gray!8},
  breaklines=true
}

\title{Labwork 2: Get to know your GPU}
\author{Yanis Denizot -- 2640051}
\date{\today}

\begin{document}
\maketitle

\section{Objective}
The goal of this labwork is to use \texttt{numba.cuda} to get basic information
about the GPU: its name, its number of multiprocessors and cores, and its memory size.


\section{Environment}
My laptop only has an Intel Iris Xe graphics card, and CUDA only works with NVIDIA GPUs.
So I ran the code on Google Colab (Linux), which gives access to an NVIDIA Tesla T4 GPU.


\section{Method}
First, \texttt{cuda.detect()} lists the available GPUs. Then \texttt{cuda.select\_device(0)}
selects the first GPU and returns a device object. From this object I read the \texttt{id} and
\texttt{name} attributes, the number of streaming multiprocessors (SM) with
\texttt{MULTIPROCESSOR\_COUNT}, and the compute capability with \texttt{compute\_capability}.
For the memory, \texttt{cuda.current\_context().get\_memory\_info()} returns the free and total
memory in bytes, so I divide by $1024^3$ to get the value in GB.

Numba does not give the number of cores directly. A GPU is divided into SMs, and each SM
contains a fixed number of cores. This number depends on the GPU architecture, which is given
by the compute capability. Compute capability 7.5 corresponds to the Turing architecture,
which has 64 CUDA cores per SM (NVIDIA documentation). So I computed:
\[
\text{total cores} = \text{number of SM} \times \text{cores per SM} = 40 \times 64 = 2560
\]
This matches the official Tesla T4 datasheet, which gives 2560 CUDA cores.
The value 64 is only valid for this GPU: on a GPU with another compute capability,
it has to be changed.


\begin{lstlisting}[language=Python]
from numba import cuda

cuda.detect()

gpu = cuda.select_device(0)
print(gpu.id)
print(gpu.name)

sm_count = gpu.MULTIPROCESSOR_COUNT
print(sm_count)

cc = gpu.compute_capability
print(cc)

# Tesla T4 : compute capability 7.5 = Turing architecture = 64 cores per SM
cores_per_sm = 64
cores = sm_count * cores_per_sm
print(cores)

free, total = cuda.current_context().get_memory_info()
print(total / 1024**3)
print(free / 1024**3)
\end{lstlisting}

\section{Results}

\begin{table}[h]
\centering
\begin{tabular}{ll}
\toprule
\textbf{Property} & \textbf{Value} \\
\midrule
Device ID          & 0 \\
Device name        & Tesla T4 \\
Compute capability & 7.5 \\
Multiprocessors (SM) & 40 \\
Cores per SM       & 64 \\
Total cores        & 2560 \\
Total memory       & 14.56 GB \\
Free memory        & 14.46 GB \\
\bottomrule
\end{tabular}
\caption{Information about the GPU}
\end{table}

Output of \texttt{cuda.detect()}:

\begin{lstlisting}
Found 1 CUDA devices
id 0             Tesla T4                              [SUPPORTED]
                      Compute Capability: 7.5
                           PCI Device ID: 4
                              PCI Bus ID: 0
                                    UUID: GPU-60411518-bba0-87d6-6592-1682f007169f
                                Watchdog: Disabled
             FP32/FP64 Performance Ratio: 32
Summary:
        1/1 devices are supported
\end{lstlisting}

\section{Conclusion}
The Tesla T4 has 40 multiprocessors, 2560 cores and about 14.5 GB of memory. This is many
more cores than a normal CPU, which is why a GPU can be much faster for tasks that can be
split into many small independent parts, like processing all the pixels of an image.


\end{document}
